\section{Conclusion}
\label{sec:conclusion}

We asked whether prompting a frontier LLM with a user's exported listening history yields recommendations as good as classical Spotify-style baselines. On $N{=}24$ \lastfm users with real $30$-day held-out windows, ranking the same $20$-track pool with seven methods, \claude with history achieves the highest \ndcg ($0.375$) and \mrr ($0.458$), beats popularity decisively ($+0.207$ \ndcg, $p{=}0.028$), and matches \cbf and \itemknn within the resolution of an $N{=}24$ paired test. The result generalizes across LLM families: \gpt is statistically indistinguishable on accuracy, while $3\times$ faster. \claude does not exploit a popularity shortcut (\arp $58.6 < $ \cbf $59.7$), and the personalization gain over a no-history ablation is large and highly significant ($+0.29$ \ndcg, $d{=}0.86$, $p{=}0.002$), establishing that the LLM is genuinely conditioning on the listening log.

\para{Key takeaway.} For ranking, off-the-shelf LLM prompting is genuinely competitive with the kinds of algorithms that historically powered platform-style music recommenders, at a per-user cost of cents and a latency of seconds.

\para{Future work.} Three extensions would tighten and broaden this result. First, scaling to $N \ge 100$ users for power on small effects, with the position-bias bootstrapping of \citet{hou2024llm} to reduce variance on the n.s.\ comparisons. Second, replacing \spotcat with the live Spotify Web API to measure hallucination on a complete catalog, and post-2024 listening data to remove pre-training contamination risk. Third, a subjective user study in the \citet{boadana2025llm} style on the same users, to test whether LLMs win on \emph{satisfaction} even at parity on objective accuracy. Tool-using and generative-retrieval LLM architectures \citep{doh2025talkplay,palumbo2025text2tracks} are natural points of comparison once the prompt-only baseline established here is in place.
